\appendixsection{Фрагменты исходного кода}

В приложении приведены ключевые фрагменты разработанного модуля интеграции.
Рисунок~\ref{src:native} --- объявление нативных методов доступа к библиотеке Vortex;
рисунок~\ref{src:readbatch} --- метод чтения очередного пакета строк с экспортом в
раскладку Apache Arrow без копирования; рисунок~\ref{src:create} --- фрагмент
параметров создания таблицы near-real-time-сценария.

\begin{figure}
\begin{lstlisting}[language=Java]
/** Нативные методы доступа к файлам формата Vortex. */
public final class NativeFileMethods {
    static { NativeLoader.loadJni(); }
    private NativeFileMethods() {}

    public static native List<String> listVortexFiles(String uri, Map<String,String> options);
    public static native void   delete(String[] uris, Map<String,String> options);
    public static native long   open(String uri, Map<String,String> options);
    public static native long   rowCount(long pointer);
    public static native long   dtype(long pointer);
    public static native void   close(long pointer);
    public static native long   scan(long pointer, List<String> columns,
                                     byte[] predicateProto,
                                     long[] rowRange, long[] rowIndices);
}
\end{lstlisting}
\caption{Объявление нативных методов (\texttt{NativeFileMethods})}
\label{src:native}
\end{figure}

\begin{figure}
\begin{lstlisting}[language=Java]
@Override
public FileRecordIterator<InternalRow> readBatch() throws IOException {
    if (!arrayIterator.hasNext()) return null;
    releaseCurrentArray();
    Array array = arrayIterator.next();          // нативный пакет
    this.currentArray = array;
    VectorSchemaRoot vsr = array.exportToArrow(allocator, reuse); // Arrow C Data Interface, zero-copy
    this.reuse = vsr;
    Iterator<InternalRow> rows = arrowBatchReader.readBatch(vsr).iterator();
    return new FileRecordIterator<InternalRow>() {
        @Override public long returnedPosition() { return returnedPosition; }
        @Override public Path filePath()         { return filePath; }
        @Override public InternalRow next() {
            if (!rows.hasNext()) return null;
            returnedPosition = positionIterator.next();
            return rows.next();
        }
        @Override public void releaseBatch() { releaseCurrentArray(); }
    };
}
\end{lstlisting}
\caption{Чтение очередного пакета строк с экспортом в Apache Arrow (\texttt{VortexRecordsReader})}
\label{src:readbatch}
\end{figure}

\begin{figure}
\begin{lstlisting}[language=SQL]
CREATE TABLE bench_table ( /* 25 столбцов */, PRIMARY KEY (id) NOT ENFORCED) WITH (
  'file.format' = 'vortex',           -- варьируется: parquet | orc | vortex
  'bucket' = '...', 'bucket-key' = 'id',
  'sequence.field' = 'seq', 'rowkind.field' = 'op',
  'changelog-producer' = 'none', 'deletion-vectors.enabled' = 'false',
  'snapshot.num-retained.min' = '500', 'snapshot.num-retained.max' = '10000',
  'snapshot.time-retained' = '24h', 'snapshot.expire.execution-mode' = 'async',
  'manifest.merge-min-count' = '1000', 'consumer.expiration-time' = '24h',
  'write-buffer-size' = '1GB',
  'metadata.stats-mode' = 'full',     -- полные статистики (без усечения строк)
  'write.batch-size' = '65536'        -- крупный пакет: меньше пересечений границы JNI
);
\end{lstlisting}
\caption{Фрагмент параметров создания таблицы near-real-time-сценария}
\label{src:create}
\end{figure}
